%% kapa: exam=1 round=30 year=2019 period=2 form=A subject=accounting subject_ko=회계학 seq=25 number=65
%% source: Q-Net 감정평가사 기출문제 (cst00309) artlSeq=5205205 file=extracted/1008/5205205/2019년 제30회 감정평가사 1차시험 문제지/2019년 제30회 감정평가사 1차시험 문제지/2019년도 제30회 감정평가사 1차 2교시 A형.hwp
%% conversion: kordoc 4.24.51 -> tools/gen1.py (auto); review=proofread-sonnet; flags=none
\begin{question}{65}
(주)감평은 20x1년 12월 31일 주거래은행으로부터 당좌예금잔액증명서상 잔액이 ￦7,810,000이라는 통지를 받았으나, 회사의 12월 31일 현재 총계정원장상 당좌예금 잔액과 불일치하였다. (주)감평이 이러한 불일치의 원인을 조사한 결과 다음과 같은 사항을 발견하였다. 이들 자료를 활용하여 (주)감평의 수정 전 당좌예금계정 잔액(A)과 수정 후 재무상태표에 당좌예금으로 계상할 금액(B)은?
\begin{framedtext}
○ (주)감평이 발행하고 인출 기록한 수표 ￦2,100,000이 은행에서 아직 지급되지 않았다. \par
○ 매출거래처로부터 받아 예금한 수표 ￦1,500,000이 부도 처리되었으나, (주)감평의 장부에 기록되지 않았다. \par
○ 주거래은행에 추심의뢰한 받을어음 ￦500,000이 (주)감평의 당좌예금 계좌로 입금 처리되었으나, 통보받지 못하였다. \par
○ 지난 달 주거래은행에 현금 ￦190,000을 당좌예입하면서 회계직원의 실수로 장부상 ￦910,000으로 잘못 기장된 것이 확인되었다. \par
\end{framedtext}
\begin{choices}
  \item A: ￦5,990,000 B: ￦5,210,000
  \item A: ￦5,990,000 B: ￦5,710,000
  \item A: ￦7,430,000 B: ￦5,710,000
  \item A: ￦7,430,000 B: ￦6,430,000
  \item A: ￦9,530,000 B: ￦7,310,000
\end{choices}
\end{question}
\answer{3}
